% ==================================================
% Telescoping Sums
% Discrete Calculus
% ==================================================

\section{Telescoping Sums}
\label{sec:telescoping-sums}

One of the most useful patterns in summation is the phenomenon of
\emph{telescoping}. In a telescoping sum, consecutive terms cancel with
one another, leaving only a small number of terms from the beginning
and end of the expression.

\section*{A Simple Example}

Consider the sum
\[
(2-1) + (3-2) + (4-3) + (5-4).
\]
Expanding gives
\[
2 - 1 + 3 - 2 + 4 - 3 + 5 - 4.
\]
Most terms cancel in adjacent pairs:
\[
\cancel{2} - 1 + \cancel{3} - \cancel{2} + \cancel{4} - \cancel{3} + 5 - \cancel{4}.
\]
After cancellation, only the first and last terms remain: \(5 - 1\).
Thus
\[
(2-1) + (3-2) + (4-3) + (5-4) = 5 - 1.
\]

\begin{tcolorbox}[colback=defbox, colframe=defborder, arc=2pt,
  left=6pt, right=6pt, top=4pt, bottom=4pt,
  title={\small\textbf{Definition (Telescoping Sum)}},
  fonttitle=\small\bfseries]
A sum is called a \emph{telescoping sum} if consecutive terms cancel when
the sum is expanded, leaving only a few remaining terms — typically
the first and last.
\end{tcolorbox}
\label{def:telescoping-sum}

\section*{General Form}

The general telescoping structure is
\[
\sum_{k=a}^{b-1} (f(k+1) - f(k)).
\]
Writing out the terms explicitly:
\[
\begin{aligned}
& (f(a+1)-f(a)) \\
& + (f(a+2)-f(a+1)) \\
& + (f(a+3)-f(a+2)) \\
& \quad \vdots \\
& + (f(b)-f(b-1)).
\end{aligned}
\]
All intermediate terms cancel in pairs, leaving only \(f(b) - f(a)\).

\begin{remark*}[Why ``telescoping'']
The name comes from an old collapsible telescope: each section slides
into the next. The intermediate terms in the sum ``slide together''
and vanish, collapsing the entire sum down to its endpoints.
\end{remark*}
